%% =====================================================================
%%  B3-T-08-02 -- what B3 contributes to "The Engagement Decision"
%%  -------------------------------------------------------------------
%%  LAYER A: APPEND-ONLY. Written once, then left alone. Material found
%%  only here sits in the `unique` slot and is protected by rule R10.
%% =====================================================================
%% @kind: source
%% @id: B3-T-08-02
%% @book: B3
%% @topic: T-08-02
%% @locator: Law 36, pp. 300--308; Law 20; Law 22
%% @status: distilled
%% @unique: yes
%% @integrated: yes
%% @updated: 2026-09-19

\dossier{B3}{T-08-02}{\bookshort{B3}}{Law 36, pp. 300--308; Law 20; Law 22}

\begin{distilled}
  \item Acknowledging a petty problem gives it existence and credibility; the more attention paid to an enemy, the stronger they become.
  \item A small mistake often becomes worse and more visible when you try to fix it, so leaving things alone is sometimes best.
  \item Showing contempt for something you cannot have is presented as raising your apparent standing: the less interest revealed, the more superior you seem.
  \item Not committing to anyone is recommended, with independence described as the position from which others can be played against one another.
  \item Surrender is treated as a tool where you are weaker: it buys time to recover and to wait for the other party's power to wane, and denies them the satisfaction of a victory.
\end{distilled}

\begin{unique}
  \item Non-response as an active instrument --- attention withheld is described as costing the other party more than a reply would.
  \item Surrender reframed as a chosen timing decision rather than as a defeat.
\end{unique}
